\DocumentMetadata{
  pdfversion=2.0,
  pdfstandard=ua-2,
  lang=en-US,
  tagging=on,
  tagging-setup={math/setup=mathml-SE,tabsorder=structure}
}
\documentclass[11pt,letterpaper]{article}
\usepackage[margin=0.68in,includefoot]{geometry}
\usepackage{newcomputermodern}
\usepackage{amsmath,amscd,mathtools}
\usepackage{tikz,adjustbox}
\providecommand{\square}{\mdlgwhtsquare}
\usepackage[hidelinks]{hyperref}
\hypersetup{
  pdftitle={Analysis First-Year Exam, August 2015, Part 1},
  pdfauthor={University of Florida Department of Mathematics},
  pdfsubject={Graduate mathematics examination},
  pdfdisplaydoctitle=true,
  bookmarksopen=true
}
\setcounter{secnumdepth}{0}
\setlength{\parindent}{0pt}
\setlength{\parskip}{0.28em}
\setlength{\emergencystretch}{3em}
\setlength{\leftmargini}{2.15em}
\setlength{\leftmarginii}{2.65em}
\setlength{\leftmarginiii}{3em}
\pagestyle{empty}
\raggedbottom
\allowdisplaybreaks
\NewTaggingSocketPlug{block/list/label}{examproblem}{%
  \tagstructbegin{tag=Lbl}\tagstructend%
  \tagstructbegin{tag=\LBody}%
  \tagstructbegin{tag=\ProblemHeadingTag,title-o={\ProblemHeadingText}}%
  \tagmcbegin{tag=\ProblemHeadingTag,actualtext={\ProblemHeadingText}}%
  #2%
  \tagmcend%
  \tagstructend%
}
\newenvironment{examproblems}[2]{%
  \def\ProblemHeadingTag{#2}%
  \AssignTaggingSocketPlug{block/list/label}{examproblem}%
  \begin{enumerate}%
  \setcounter{enumi}{#1}%
  \renewcommand{\labelenumi}{\textbf{\theenumi.}}%
  \setlength{\topsep}{0.35em}%
  \setlength{\partopsep}{0pt}%
  \setlength{\itemsep}{0.48em}%
  \setlength{\parsep}{0.18em}%
}{\end{enumerate}}
\newenvironment{examsubparts}{%
  \AssignTaggingSocketPlug{block/list/label}{default}%
  \begin{enumerate}%
  \setlength{\topsep}{0.18em}%
  \setlength{\partopsep}{0pt}%
  \setlength{\itemsep}{0.16em}%
  \setlength{\parsep}{0.1em}%
}{\end{enumerate}}
\newenvironment{examinstructions}{%
  \par\begingroup\itshape
}{%
  \par\endgroup\vspace{0.18em}
}
\makeatletter
\renewcommand\subsection{\@startsection{subsection}{2}{0pt}%
  {-1.25ex plus -.25ex minus -.1ex}%
  {0.55ex plus .1ex}%
  {\normalfont\large\bfseries}}
\renewcommand{\@maketitle}{%
  \newpage\null\vskip -1em%
  {\footnotesize\bfseries
    \href{https://www.ufl.edu/}{University of Florida}, \href{https://math.ufl.edu/}{Department of Mathematics}\par}%
  \vskip -0.5em%
  \tagstructbegin{tag=H1,title={Analysis First-Year Exam, August 2015, Part 1}}%
  \tagpdfparaOff%
  \tagmcbegin{tag=H1}%
    {\LARGE\bfseries\href{https://gma.math.ufl.edu/exams/analysis-first-year/}{Analysis First-Year Exam}, August 2015, Part 1\par}%
  \tagmcend%
  \tagpdfparaOn%
  \tagstructend%
  \vskip 0.25em%
  \tagmcbegin{artifact}%
    \hrule height 1pt%
  \tagmcend%
  \vskip 1em%
}
\makeatother
\title{Analysis First-Year Exam, August 2015, Part 1}
\author{}
\date{}
\begin{document}
\maketitle
\thispagestyle{empty}
\begin{examinstructions}
Do exactly 4 problems. Work must be presented in a neat and logical fashion in order to receive credit. Do not leave any gaps. When a theorem is used in a proof it must be precisely stated.
\end{examinstructions}
\begin{examproblems}{0}{H2}
\def\ProblemHeadingText{Problem 1.}
\item
Let~\(E\) be a dense subset of a metric space~\(X\) and let~\(U\) be an open subset of~\(X\). Prove that \(U \subset \overline{E \cap U}\). (Here \(\overline{E \cap U}\) denotes the closure of \(E \cap U\).)
\def\ProblemHeadingText{Problem 2.}
\item
Let~\(X\) be a metric space and \(K \subset X\) a compact subset. Prove that~\(K\) is closed.
\def\ProblemHeadingText{Problem 3.}
\item
Let~\(X\) be a compact metric space and \(\mathcal F\) a collection of closed subsets of~\(X\) with the following property: for all finite subcollections \(\mathcal G \subset \mathcal F\), the intersection \(\bigcap_{F \in \mathcal G} F\) is nonempty. Prove that \(\bigcap_{F \in \mathcal F} F\) is nonempty.
\def\ProblemHeadingText{Problem 4.}
\item
Let \(X,Y\) be metric spaces and \(f:X\to Y\) a continuous function. Prove that if~\(X\) is connected, then its image \(f(X)\) is connected.
\def\ProblemHeadingText{Problem 5.}
\item
Let \((a_n)_{n\in\mathbb N}\) be a sequence of real numbers and~\(\sigma\) a permutation of \(\mathbb N\). For each of the following provide either a proof or a counterexample.
\begin{examsubparts}
\item[\textnormal{a)}]
If the sequence \((a_n)_{n\in\mathbb N}\) converges then so does the rearrangement \((a_{\sigma(n)})_{n\in\mathbb N}\).
\item[\textnormal{b)}]
If the series \(\sum_{n\in\mathbb N}a_n\) converges then so does the rearrangement \(\sum_{n\in\mathbb N}a_{\sigma(n)}\).
\end{examsubparts}
\def\ProblemHeadingText{Problem 6.}
\item
Let \(f(x)=x^2\sin\frac{1}{x}\) for \(x\in\mathbb R\setminus\{0\}\) and \(f(0)=0\). Prove that~\(f\) is differentiable at every \(x\in\mathbb R\). Is \(f'\) continuous at \(x=0\)? (Prove your claim.)
\end{examproblems}
\end{document}
